
% In this section, we briefly compare our results to selected previous literature
% on the BCLT.

% Standard BCLTs which approximate the posterior are not actaully so different
A standard approach to BCLTs emphasizes the large-data distributional limit of
the Bayesian posterior, e.g., by showing that the total variation distance
between the Bayesian posterior and a normal approximation convergences to zero
in probability.  For example, this is the approach taken by Chapter 10 of
\citet{van:2000:asymptotic}, Theorem 1.4.2 of \citet{ghosh:2003:bayesian}, and
\citet{lehman:1983:pointestimation}, on which our proof is principally based.
Results such as our \thmref{bayes_clt_main}, which emphasize the frequentist
properties of posterior quantities, may seem to be of a different character than
distributional approximations to the posterior itself.  However, the approach to
proving convergence in total variation is essentially to prove that posterior
expectations of bounded functions converge uniformly.  Though our assumption of
the differentiability $\psi(\theta, \xvec)$ proscribes the direct use of our
$\thmref{bayes_clt_main}$ to prove convergence in total variation distance, the proof
technique is the same, differing only in whether $\psi(\theta, \xvec)$ is Taylor
expanded in $\theta$ near $\thetahat$.  Indeed, one might argue that many BCLTs
consist of two distinct components: first, to show that only a vanishing ball
around $\thetahat$ matters for the posterior, and, second, to exploit
local properties of the posterior near $\thetahat$ to prove whatever result
is desired.  The first step is, in a sense, generic, and extra assumptions
can be made (or not) in a region of $\thetahat$ to achieve the desired
theoretical aim.  For example, total variation distance forgoes the
expansion of $\psi(\theta, \xvec)$, and a version of our \thmref{bayes_clt_main}
expanding to higher orders requires nothing more than gather the next order
terms in a the local Taylor expansion of posterior quantities.

% The global uniqueness is a little weird.
As mentioned above, half of a proof of a BCLT typically involves proving that
the behavior of the posterior far from $\thetahat$ is negligible.  To this end,
some assumptions must be made to ensure that the posterior is well behaved
outside a ball centered at $\thetahat$.  For example, one must ensure that there
are no other posterior modes very nearly as high as that at $\thetahat$.  In the
present work we follow \citet{ghosh:2003:bayesian, lehman:1983:pointestimation}
and manage the posterior far from $\thetahat$ with the final condition of
\assuref{bayes_clt}, namely that $\thetatrue$ is a strict optimum of the
empirical likelihood with probability approaching one.  This assumption may seem
somewhat artificial, since it is stated in terms of the data, rather than of any
limiting population quantities.  Recent authors, notably
\citet{kleijn:2006:misspecification}, express a similar requirement in terms of
the existence of a consistent test for $\theta = \thetatrue$, which is more
general and arguably cleaner, but introduces extra complications and  perhaps
obscure somewhat the role of the assumption in the proof.  Suffice to say, the
uniqueness of the optimum in misspecified BCLTs is a serious issue, but a
complicated one that is beyond the scope of the present paper, and we
essentially assume the problem away.  Again, see
\citet{kleijn:2006:misspecification} for a more thorough discussion.
